% IEEE appendix input. Compile inside the main document.
\section{Historical Circuit Variants}
\label{sec:historical}
Two early receiver families share a resistively loaded differential pair.
The load experiment exposes the two drain voltages directly. The
direct-coupled clock receiver adds two CMOS inverters at the positive
output (Fig.~\ref{fig:hist-rx}). In the selected receiver, a coupling
capacitor and a feedback resistor set the first inverter's input bias
(Fig.~\ref{fig:aprx}). This distinction matters when comparing the
measured loading and clock waveforms.

\begin{figure*}[!t]
\centering\resizebox{.90\textwidth}{!}{%
\begin{circuitikz}[american,font=\small]
% Direct CML: positive input drives the negative output drain.
\node[nmos] (hip) at (0,1.5) {}; \node[nmos,xscale=-1] (hin) at (2.8,1.5) {};
\node[nmos] (ht) at (1.4,-.15) {};
\draw (hip.S) -- (0,.66) -- (2.8,.66) -- (hin.S);
\draw (1.4,.66) node[circ]{} -- (ht.D);
\draw (ht.S) -- ++(0,-.18) node[below]{VSS};
\draw (ht.G) -- ++(-.15,0) node[left]{BN};
\draw (hip.G) -- ++(-.12,0) node[left]{INP};
\draw (hin.G) -- ++(.12,0) node[right]{INN};
\draw (hip.D) -- (0,2.25) node[circ]{} to[R] (0,3.9);
\draw (hin.D) -- (2.8,2.25) node[circ]{} to[R] (2.8,3.9);
\draw (0,3.9) -- (2.8,3.9) node[midway,above]{VDD};
\draw (0,2.25) -- ++(-.6,0) node[left]{ON};
\draw (2.8,2.25) -- ++(.55,0) node[right]{OP};
\node[align=center] at (1.4,-1.88) {(a) Shared CML front end};
% Repeated labels refer to the same OP drain. Do not connect OP through a capacitor.
\begin{scope}[xshift=6.4cm]
\node[american not port] (a) at (1.1,2.45) {};
\node[american not port] (b) at (3.9,2.45) {};
\draw (-.63,2.45) node[left]{OP} -- (a.in);
\draw (a.out) -- node[above]{$\overline{\mathrm{CLK}}$} (b.in);
\draw (b.out) -- ++(.3,0) node[right]{CLK};
\node[align=center] at (2.05,1.05) {Direct receiver extension\\OP drives the first inverter gate};
\node[align=center] at (2.00,-1.88) {(b) Two added CMOS inverters};
\end{scope}
\begin{scope}[xshift=13.5cm]
\node[nmos,rotate=90] (hs) at (1.3,2.75) {};
\draw (hs.D) -- ++(-.29,0) node[left]{CTRL};
\draw (hs.G) -- ++(0,-.13) node[below]{$s_{k+}$};
\draw (hs.S) -- (3.2,2.75) node[circ]{} -- ++(.43,0) node[right]{$g_k$};
\node[nmos] (hg) at (3.2,1.13) {};
\draw (hg.D) -- (3.2,2.75); \draw (hg.G) -- ++(-.1,0) node[left]{$\overline{s}_{k+}$};
\draw (hg.S) -- ++(0,-.19) node[below]{VSS};
\node[align=center] at (1.75,-1.88) {(c) Early fine-bank gate drive};
\end{scope}
\end{circuitikz}}
\caption{Archived load and oscillator variants. With $1$-k$\Omega$ loads the CML-only cell has four input fingers of either 1.5 or 3~$\mu$m per device. The direct receiver defaults to $3$-k$\Omega$ loads and adds (b). Both use a 12-finger tail with 2.5-$\mu$m fingers. In (c), each of $k=8,16,24$ uses an NMOS to pass CTRL to a supplemental sink gate and a second NMOS to discharge it. The later ring adds a parallel PMOS on the CTRL path.}
\label{fig:hist-rx}
\end{figure*}

\begin{figure*}[!t]
\centering\resizebox{.80\textwidth}{!}{%
\begin{circuitikz}[american,font=\small,
 tx/.style={draw,align=center,inner sep=3pt,minimum width=1.8cm,minimum height=.58cm}]
% Selector TG inputs are tied to named rail or mirror nodes.
\node[tx] (ub) at (1.7,3.35) {TG, UP};
\node[tx] (ur) at (1.7,2.40) {TG, $\overline{\rm UP}$};
\node[tx] (db) at (1.7,1.0) {TG, DN};
\node[tx] (dr) at (1.7,.05) {TG, $\overline{\rm DN}$};
\draw (ub.west) -- ++(-.42,0) node[left]{$b_{p,\rm cp}$};
\draw (ur.west) -- ++(-.42,0) node[left]{VDD};
\draw (db.west) -- ++(-.42,0) node[left]{$b_{n,\rm cp}$};
\draw (dr.west) -- ++(-.42,0) node[left]{VSS};
\draw (ub.east) -- (3.2,3.35) -- (3.2,2.40) -- (ur.east);
\draw (db.east) -- (3.2,1.0) -- (3.2,.05) -- (dr.east);
\node[pmos] (mu) at (5.15,2.86) {}; \node[nmos] (md) at (5.15,.53) {};
\draw (3.2,2.86) node[circ]{} -- node[above]{$g_p$} (mu.G);
\draw (3.2,.53) node[circ]{} -- node[above]{$g_n$} (md.G);
\draw (mu.S) -- ++(0,.32) node[above]{VDD}; \draw (md.S) -- ++(0,-.32) node[below]{VSS};
\draw (mu.D) -- (md.D); \draw (5.15,1.7) node[circ]{} -- ++(.5,0) node[right]{CTRL};
\begin{scope}[xshift=8.4cm]
\node[american not port] (iu) at (1.2,2.9) {};
\node[american not port] (idn) at (1.2,1.80) {};
\draw (iu.in) -- ++(-.3,0) node[left]{UP}; \draw (iu.out) -- ++(.3,0) node[right]{$\overline{\rm UP}$};
\draw (idn.in) -- ++(-.3,0) node[left]{DN}; \draw (idn.out) -- ++(.3,0) node[right]{$\overline{\rm DN}$};
\node[align=center] at (1.18,.30) {Mirror bias network\\as in Fig.~\ref{fig:apcp}};
\end{scope}
\end{circuitikz}}
\caption{Gate-switched charge-pump candidate. Each of the four TGs uses an NMOS/PMOS pair with complementary gate drives. The label gives the NMOS enable. With UP high the PMOS output gate selects the mirror voltage, and with UP low it selects VDD. DN likewise selects the NMOS mirror voltage or VSS. The TG and inverter widths are $0.5/1.5~\mu$m at $L=0.13~\mu$m. The NMOS and PMOS output widths are 4 and 12~$\mu$m at $L=1.04~\mu$m.}
\label{fig:hist-cp}
\end{figure*}

The original AFC core R0 contains 6,908 MOS instances and 145 DFFs.
Its FB counter uses the physical FB clock directly. Reset follows
$\mathrm{RST}\lor\overline{\mathrm{count\_run}}$. R1 adds a latched
FB-clock enable and holds counter reset deasserted during DRAIN,
as shown in Fig.~\ref{fig:r1afc-datapath}. The archive also records
the R2/counter9 proposal as a prepared design option. The nominal
closed-loop result in this report belongs to R1.

The simulator additionally accepts eight project-authored behavioral
modules. They represent a charge pump, precharge path, linear VCO,
divide-by-32 counter, PFD, search controller, handoff controller and
AFC unit-test source. Their equations and event handlers provide stimuli
and functional substitutes in staged benches. The transistor counts in
the selected design apply to the MOS implementations.

The companion \texttt{module-index.md} maps all 32 distinct Spectre
subcircuit names to these drawings and identifies the eight Verilog-A
modules by name.
